\documentclass[10pt]{article}
\usepackage[margin=20mm]{geometry}
\usepackage{amsmath,amssymb,mathtools,fontspec,unicode-math,xurl,hyperref}
\IfFontExistsTF{DejaVu Serif}{\setmainfont{DejaVu Serif}}{\setmainfont{Latin Modern Roman}}
\setmathfont{latinmodern-math.otf}
\setlength{\parindent}{0pt}\setlength{\parskip}{6pt}\sloppy
\title{Formal initial review: M25-0074
The Smale conjecture and min-max theory}
\author{Scientific review performed in this conversation}
\date{26 October 2026}
\begin{document}
\maketitle
\section*{Scope and initial outcome}
The acquired primary PDF is the 38-page author manuscript identified above. The journal DOI is https://doi.org/10.1007/s00222-025-01334-z. The publisher landing page identifies an article published on 4 April 2025 in Inventiones mathematicae 241, 1–34; the journal full text was not acquired or established to coincide with this manuscript. All findings here concern v3 only. PDF and printed page numbers coincide.\par
Initial outcome: Major. The explicit homotopy used in Proposition 7.2 can leave its prescribed target. A smooth relative family satisfying the normalization in the proof gives an exact counterexample. This invalidates a displayed construction in the principal proof chain. A range-preserving replacement for that local construction is proved below; the main theorem is not disproved. Four other literal statements have local scope or orientation defects, with counterevidence restricting their impact.\par
No. | Confirmed statement or step | Assessment\par
1 | Collar homotopy, p. 33, between (7.13) and (7.14) | Major; core construction\par
2 | Uniform absolute projection-area estimate (6.1), p. 22 | Minor; excess direction unnecessary in the main use\par
3 | Boundary equality in Lemma 6.4, p. 27 | Minor; missing boundary hypothesis\par
4 | All essential embeddings versus an oriented target, pp. 30–32 | Minor; excess component in Proposition 7.2\par
5 | Connected identity-orbit quotient versus two points, pp. 6–7 | Minor; incompatible target definitions\par
All 38 pages were read as extracted text, including references and the appendix. Eighteen source pages were inspected as actual images. The acquired TeX/bibliography tree was parsed completely for structure, dependencies, and execution tokens; critical passages were also read directly. No author code or author TeX was executed. Absence of computational code is not a defect.\par
The argument reviewed runs from the diffeomorphism/embedding pair reduction, through genus-one min–max and the near-minimal relative cycle, to removal of thin protrusions, collar retraction, and the final relative homotopy vanishing. The external classification and min–max theorems were checked for the roles assigned to their hypotheses, but not independently reproved in full. Detailed coverage limits appear on the last page.\par
\section*{Problem 1. The displayed collar homotopy exits the tube}
Location and claim. Proposition 7.2, pp. 32–33; specifically the unnumbered formula following (7.13) and used to define (7.14). In product coordinates, a collar \(e:T^2\times[-r,0]\to T^2\times[-r,r]\) has bottom height \(-r\), and is normalized by \(\partial_\rho e^2(y,-r)=1\). The proposed homotopy, for \(t>0\), is \[e_t(y,\rho)=\left(e^1(y,t(\rho+r)-r),
-r+\frac{e^2(y,t(\rho+r)-r)+r}{t}\right).\] The proof treats this as a path in the same collar space.\par
Complete counterexample. Put \(u=(\rho+r)/r\) and \[F(u)=u+\frac{16}{5}u^2-\frac{49}{20}u^3,\qquad
 e(y,\rho)=(y,-r+rF(u)).\] Here \(F(0)=0\), \(F'(0)=1\), \(F(1)=7/4\), and \[F'(u)=1+\frac{32}{5}u-\frac{147}{20}u^2.\] This last polynomial is concave; its endpoint values on \([0,1]\) are \(1\) and \(1/20\). Thus it is strictly positive throughout the interval. Consequently \(e\) is a smooth collar embedding, maps its lower boundary onto the lower boundary of the tube, has the required bottom derivative, and has image heights in \([-r,3r/4]\). Its top is an essential, positively oriented torus in the intended isotopy class.\par
Nevertheless, at \(\rho=0\) and \(t=2/3\), the proposed top height is \[-r+r\frac{F(2/3)}{2/3}
 =-r+\frac{92}{45}r
 =\frac{47}{45}r>r.\] It therefore does not lie in the target tube; the subsequent interpretation as an element of the original collar space is unavailable.\par
This is also a relative-family counterexample. For \(x\in[-1,1]\), let \(\lambda(x)=1-x^2\) and \[F_x(u)=u+\lambda(x)\bigl(F(u)-u\bigr).\] Then \(F'_x=(1-\lambda)+\lambda F'>0\), \(F_x(1)\le 7/4\), and the bottom conditions hold for every \(x\). At both boundary parameters \(x=\pm1\), the collar is exactly \((y,\rho)\), belonging to the standard minimal collar family. At \(x=0\), the preceding violation remains. For a higher-dimensional parameter disk use \(\lambda(x)=1-\|x\|^2\). Thus the relative boundary requirement does not exclude the example.\par
Counterevidence and impact. The formula works for a linear vertical collar, and its limit as \(t\to0\) is correct. Neither fact gives a bound for intermediate \(t\). No global inequality controlling \(e^2+r\) by \(\rho+r\) is imposed at this point; normalizing the derivative only at the bottom does not impose one. Shrinking \(r\) does not fix this scale-invariant counterexample. The defect is independent of the orientation issue in Problem 4.\par
Proposition 7.2 supplies the contractions in Proposition 7.4 (p. 34), which finishes Theorem 7.3 and Claims 4.1–4.2, hence Theorem 2.8. This is a demonstrated failure of a core construction, warranting Major under the review calibration even though a local repair exists. It is not a counterexample to the Smale conjecture.\par
\section*{Problem 1 continued: a verified local replacement}
Here is a continuous relative-family repair of the range defect. For a compact normalized family \(e_x\), set, for \(t>0\), \[A_{x,t}(y,\rho)=
 \frac{e_x^2(y,t(\rho+r)-r)+r}{t},\qquad
 M_{x,t}=\max_{(y,\rho)\in T^2\times[-r,0]}A_{x,t}(y,\rho),\] \[\ell_{x,t}=\min(1,r/M_{x,t}).\] Since the collar is contained in the original tube, \(A_{x,t}\ge0\). It is not identically zero because the normalized derivative is one; hence \(M_{x,t}>0\). These maxima are continuous in the compact family parameters. Uniform Taylor expansion at the bottom gives \[A_{x,t}\longrightarrow \rho+r,\qquad M_{x,t}\longrightarrow r,\qquad
 \ell_{x,t}\longrightarrow1\] as \(t\to0\), with the first convergence in every spatial \(C^j\) norm.\par
First deform \(e_x\) by vertical contraction about height \(-r\), using the scalar \(b_{x,s}=(1-s)+s\ell_{x,1}>0\): \[e^{\mathrm{pre}}_{x,s}(y,\rho)
 =\bigl(e_x^1(y,\rho),-r+b_{x,s}(e_x^2(y,\rho)+r)\bigr).\] This is an embedding, fixes the lower boundary, and remains in the tube. Then, with \(t\) decreasing from one to zero, use \[\widetilde e_{x,t}(y,\rho)=
 \bigl(e_x^1(y,t(\rho+r)-r),-r+\ell_{x,t}A_{x,t}(y,\rho)\bigr).\] At \(t=1\) this is the endpoint of the preliminary contraction. For \(t>0\), it is a restriction and reparametrization of an embedding, followed by a positive vertical affine rescaling; hence it remains an embedding. Its height lies in \([-r,0]\), because \(\ell_{x,t}A_{x,t}\le r\). At zero define \[\widetilde e_{x,0}(y,\rho)=(e_x^1(y,-r),\rho).\] The bottom restriction is a diffeomorphism of \(T^2\), so this limit is also an embedding. The Taylor expansion proves continuity into the smooth spatial mapping space, uniformly on the compact parameter set.\par
For a standard minimal collar, \(e^2=\rho\) and \(e^1\) is independent of \(\rho\). Then \(M_{x,t}=r\), \(\ell_{x,t}=1\), and both homotopies fix that collar. Thus the repair is relative to the same boundary family. The maximum need not be differentiable in parameters; a continuous parameterized homotopy is all that is used here.\par
Repair limits. This verifies the replacement for the displayed scaling step after the paper’s bottom-derivative normalization. It does not supply the preceding normalization argument, settle all mapping-class components, or certify the later gluing automatically. The intended orientation component must be specified as in Problem 4. No claim of full acceptance of Proposition 7.2 or of the entire theorem follows from this local repair.\par
Independent computation. Reviewer-written SymPy checks evaluated the polynomial endpoints, derivatives, concavity, relative boundary values, convex-combination identity, and \(47/45\) exactly. The actual interpreter output is retained in `independent\_checks.json`. These algebra checks support, and do not replace, the embedding and continuity proof above.\par
\section*{Problem 2. Equation (6.1) has no uniform absolute error}
Location. Section 6, p. 22, (6.1); subsequent use on p. 23 in (6.4)–(6.6). The asserted bound is \[\bigl|\operatorname{Area}(\Gamma)-\|p_\#\Gamma\|\bigr|\le \Lambda r^2\] for every smooth closed surface \(\Gamma\subset N_r(\Sigma)\), with \(\Lambda\) depending only on the ambient manifold. No area or slope bound is included in that assertion.\par
Complete counterexample. Use the standard genus-one Heegaard splitting of \(S^3\). Choose a smooth Riemannian metric which is \(dx^2+dy^2+dt^2\) on a collar of the splitting torus \(\Sigma=(\mathbb R/\mathbb Z)^2\times\{0\}\). Such a positive smooth metric extends over the two solid tori. The surface remains a Heegaard torus and has area one. Fix a sufficiently small \(r>0\), set \(a=r/2\), and take \[\Gamma_N=\{(x,y,a\sin(2\pi Nx)):x,y\in\mathbb R/\mathbb Z\},
 \qquad N\in\mathbb N.\] Each is a connected smooth embedded closed torus in the prescribed tube. Projection has degree one and \(\|p_\#\Gamma_N\|=1\), whereas \[\operatorname{Area}(\Gamma_N)=
 \int_0^1\sqrt{1+(2\pi Na\cos(2\pi Nx))^2}\,dx
 \ \ge\ 4aN.\] Therefore the absolute difference is at least \(4aN-1\), unbounded for fixed \(r\). Taking \(N>(\Lambda r^2+1)/(4a)\) violates the assertion. This checks all assumptions of the standalone estimate. It does not assert that these highly oscillatory graphs satisfy the later small-\(\mathbf F\) hypothesis of Proposition 6.1.\par
Counterevidence: the main application needs less. The use in (6.4) needs a lower area bound arising from projection degree, not a uniform upper bound for the graph’s area excess. For a round Clifford torus \(C\), Fermi coordinates give the metric \[ds^2+\cos^2(\pi/4+s)\,d\alpha^2+
       \sin^2(\pi/4+s)\,d\beta^2.\] The two-dimensional Jacobian of projection on any tangent two-plane is at most \(\sec(2r)\) for \(|s|\le r<\pi/4\). Indeed, its maximum is on the tangent plane of the parallel torus, whose area-density ratio to \(C\) is \(\cos(2s)\); planes with a normal component cannot increase that Jacobian. Thus \[\|p_\#\Gamma\|\le\sec(2r)\operatorname{Area}(\Gamma).\] When the projection has degree one, this gives precisely \[\operatorname{Area}(\Gamma)\ge\cos(2r)\operatorname{Area}(C)
 =\operatorname{Area}(C)-O(r^2).\] This proves the needed one-sided estimate in the main round-Clifford application, without a slope restriction.\par
Impact and repair limits. Replace (6.1) by the appropriate Jacobian inequality. For a general nonminimal Heegaard surface the analogous distortion is usually \(1+O(r)\), so the full generality of Proposition 6.1 still needs an adjusted smallness argument. No counterexample to that proposition was established. The confirmed false part of (6.1) is an overstrong auxiliary estimate whose unneeded direction can be removed in the main application; it is classified Minor.\par
\section*{Problem 3. Lemma 6.4 lacks a boundary membership hypothesis}
Location and claim. Page 27, Lemma 6.4 asks for a map \(g:X\to\mathcal T(\Sigma,M)\) agreeing with the oriented image of \(f(x)\) on \(\partial X\), but assumes only that \(f\) is sufficiently \(\mathbf F\)-close to that isometric orbit.\par
Counterexample with all assumptions. Let \(X=[0,1]\), \(M=L(3,1)\) with its round metric, and let \(\Sigma=C\) be a Clifford torus. Take \(f\) constant, equal to a parametrization of the parallel torus \(C_s\), where \(s\ne0\) is arbitrarily small. This is an embedding in the required global component and \(\mathbf F(C_s,C)\to0\). It satisfies any prescribed positive closeness tolerance by choosing \(s\) sufficiently small. However, \(\operatorname{Area}(C_s)=\cos(2s)\operatorname{Area}(C)<\operatorname{Area}(C)\). Every member of the isometric orbit \(\mathcal T(\Sigma,M)\) has area \(\operatorname{Area}(C)\). Thus \(C_s\) is not in that orbit, and the required equality \(g(0)=C_s\) is impossible.\par
Repair and counterevidence. Add \(f(\partial X)\subset\mathcal T(\Sigma,M)\), interpreting the image with the chosen orientation. Set \(g=f\) on the boundary before performing the relative small-cell extension. Proposition 6.5 on the same page already includes this extra boundary hypothesis, so its stated application is compatible with the intended version. The contradiction is to the lemma as written, not to the main theorem. This is a Minor scope defect. A complete independent construction of all the parameter-dependent local extensions was not undertaken.\par
\section*{Problem 4. Proposition 7.2 includes an opposite orientation component}
Location and claim. Page 30 defines \(\operatorname{Emb}_0(T^2,N_r(C_q))\) using the global component in \(M\) and essential images in the tube. The target \(\mathcal T_{\min}(q)\) fixes the orientation of \(q\). Proposition 7.2, p. 32, includes \(k=0\) and asserts surjectivity onto all path components of this domain.\par
Counterexample. Take \(M=L(3,1)=S^3/\mathbb Z_3\), with central deck action, \(q=f_0\) the standard Clifford embedding, and \(\sigma(z,w)=(w,z)\). This matrix lies in the connected group \(U(2)\) and commutes with the deck action; a path in \(U(2)\) descends to a path of isometries from the identity to \(\sigma\). Hence \(E=\sigma\circ f_0\) lies in the global component \(\operatorname{Emb}_0(T^2,M)\). Its image is exactly \(C_q\), hence essential in every small tube. Swapping the two angular coordinates has determinant \(-1\), so \(E\) represents \(-[C_q]\) in \(H_2(N_r(C_q);\mathbb Z)\cong\mathbb Z\). Every element of \(\mathcal T_{\min}(q)\) represents \(+[C_q]\). Homotopies through maps into the tube preserve this homology class; there can be no such path from \(E\) to the claimed target.\par
Impact, counterevidence, and repair. Restrict the domain to projection degree \(+1\), with compatible oriented parametrizations. The proof invokes an oriented Goeritz group and cannot use its triviality for an isometry swapping the handlebodies. The preceding construction of the near-minimal family uses oriented \(\mathbf F\)-closeness, which supports the intended restricted version. A failure of the final theorem has not been shown; preservation of this restriction through every extension still requires an explicit trace. This is classified Minor, as an extra-component defect outside the intended oriented application, separately from the core range failure of Problem 1.\par
\section*{Problem 5. The defined orbit cannot be a two-point space}
Location. Page 6 defines \(\overline{\mathcal T}_{\min}=\operatorname{isom}(M)/\operatorname{isom}^+(M,T)\), where \(\operatorname{isom}(M)\) is explicitly the identity component. Proposition 2.3(3), p. 7, assigns cardinality two to this space for \(q\ne1,\ p>2\). The next paragraph describes the connected space \(\mathcal T_{\min}\cong\operatorname{isom}(M)\) as a \(T^2\)-bundle over it.\par
Verification. The continuous image of a connected group under its quotient map is connected. It cannot consist of two distinct points. Concretely, for \(L(5,2)\), the identity isometry group is the two-torus of coordinate rotations. Every such rotation preserves the standard Clifford torus and its orientation, so the stabilizer \(\operatorname{isom}^+(M,T)\) is the whole group and the quotient is one point. This also excludes a surjective bundle from the connected \(\mathcal T_{\min}\) to a two-point base.\par
Counterevidence and repair. The space of both formal orientations of the single unoriented Clifford torus does have two points. That is a different space from the defined identity-component orbit in this case. Use a one-point base for this orbit, or explicitly define a larger oriented moduli space and separate its two components from the chosen parametrized component. The main relative argument concerns the component of \(f_0\); no false main conclusion follows merely from the extra formal orientation. This definitional incompatibility is Minor. A repaired convention should be used consistently in Sections 5–7.\par
\section*{Positive checks and rejected suspicions}
Relative Hurewicz, p. 13. The general relative Hurewicz theorem uses coinvariants (the quotient by relations \(\gamma a-a\), with normal closure in degree two), not an arbitrary orbit set. The paper’s orbit wording alone is not a sufficient general argument. In this particular application, however, there is a direct justification for trivial action. In the pair \((G,H)=(\operatorname{diff}(M),\operatorname{isom}(M))\), take a loop \(\gamma\) in \(H\) and a relative representative \(f\). The homotopy \(F_t(x)=\gamma(t)f(x)\) moves the basepoint along \(\gamma\), keeps the relative boundary in \(H\), and starts and ends at \(f\). Thus the basepoint action is trivial. Transfer through the relative isomorphism (2.11) gives the same result for \((\mathcal T,\mathcal T_{\min})\). The coinvariant quotient is then the original relative group. With the inductive connectivity hypothesis, the desired Hurewicz step is justified. This is recorded as a resolved point, not a confirmed failure of the application.\par
Genus and multiplicity. The genus-one sweepouts are separating. Positive Ricci curvature rules out stable two-sided closed minimal surfaces, so the stated one-sided alternative with stable connected double cover is excluded in the width contradiction. A minimal sphere in a nontrivial lens space lifts to disjoint minimal spheres in \(S^3\), contradicting Frankel. A connected lifted minimal torus has Euler characteristic zero and is covered by the stated Brendle classification. These checks support the contradiction in (4.12)–(4.16).\par
Relative chains. The signs in (5.17)–(5.18) are consistent: the bottom face cancels, and the product boundary rule leaves \(\partial I^k\times\{0\}\). The separate \(k=1\) path argument avoids replacing reduced \(H_0\) by an inappropriate unreduced statement. No computational code is needed for these topological arguments.\par
\section*{Coverage, limits, and disposition}
Range | Work actually performed | Remaining limit\par
pp. 1–10 | Definitions, component/orientation choices, stabilizer and fibration reduction, theorem scope | External handlebody contractibility not reproved\par
pp. 11–16 | Min–max hypotheses, multiplicities, genus, width contradiction, relative Hurewicz | Full analytic min–max machinery not reconstructed\par
pp. 17–21 | Near-minimal separation, cubical signs, relative cycle, \(k=1\) case | Selection from the particular minimizing sequence remains an external-argument check\par
pp. 22–29 | Projection estimates, cap surgery, disk-family contraction, boundary extension | Full parameterized surgery compatibility not independently proved\par
pp. 30–35 | Collar spaces, exact homotopy counterexample and local repair, orientation, gluing, appendix | CW/cofibration realization of all infinite-dimensional pairs not fully verified\par
pp. 36–38 | References and version/dependency identities | No exhaustive reproof of cited literature\par
Actual source images inspected: 1, 6, 7, 9, 10, 12, 13, 15, 18, 19, 22, 23, 27, 30, 31, 32, 33, 34. Direct author-source reading included lines 791–885, 1281–1410, and the critical scaling lines 1356–1415. Complete programmatic source parsing is recorded separately and is not counted as scientific reading. Obvious typesetting slips and bibliographic mistakes alone were excluded from findings.\par
Decision field | Current state\par
Judgment | Major initial-review outcome for the bound author version; main theorem not refuted\par
Repair status | Range-preserving collar replacement proved locally; projection lower bound proved for the round Clifford application; boundary and orientation restrictions specified\par
Adoption | No scientific acceptance or final-archive action; no author contact\par
Remaining | Integrate the corrected constructions and conventions; finish the external selection, parameter-family, and pair-topology checks above\par
\section*{Supporting checks}
61955c7fba314a0dfc74592670d7d914e58
19434266e1f8589b8ddc8a783d1c0\par
The exact full hash, size 411,243 bytes, and all six acquired-file hashes are recorded in the source manifest and provenance. The acquired `main.tex` is 114,159 bytes and was not compiled or executed. Only this reviewer-written report was compiled, using the existing, authorized Tectonic 0.17.0 in untrusted mode without shell escape. The cache inventory is a superset, not a per-file access trace. Every final report page is checked as an image and recorded separately.\par
Version checks used the official arXiv record (https://arxiv.org/abs/2310.05756) and the publisher DOI page. A targeted correction search did not establish an applicable correction; this is not proof that none exists. The Hurewicz formulation was checked in A. Hatcher, Algebraic Topology, Theorem 4.37, pp. 370–371, https://pi.math.cornell.edu/\textasciitilde{}hatcher/AT/AT.pdf. An attempted web screenshot of that external PDF failed; its text was read, and no external page-image inspection is claimed.\par
All scientific review, independent calculations, compilation, and page inspection were performed by this conversation. `human\_reviews=[]`; `scientific\_decision=pending`.\par
\end{document}
